\section{Methods}

\subsection{Dataset}

The analysis uses the MCF10A single-cell expression dataset deposited in
ArrayExpress under accession E-MTAB-9861 \citep{jeffery2022cenpa}.
The data were generated by Jeffery and Almouzni to profile the transcriptional
consequences of CENP-A overexpression in a p53 wild-type or dominant-negative
background in MCF10A non-transformed human mammary epithelial cells.
The count matrix contains 9,201 cells and 40,295 genes.
Cell metadata include \texttt{sample\_name}, \texttt{p53\_status},
and \texttt{CENPA\_status}, which define six joint perturbation conditions.
Dataset integrity is verified programmatically before downstream analysis.

\subsection{Preprocessing}

The pipeline reads the H5AD structure directly with \texttt{h5py}, avoiding a
hard dependency on Scanpy while following standard preprocessing logic
\citep{scanpy_preprocessing}.
Library sizes are normalized to 10,000 counts per cell, transformed by
$\log(1+x)$, filtered by minimum gene detection (minimum 10 cells), and ranked
by log-expression variance.
The highly variable gene (HVG) set contains 2,000 genes plus anchor genes:
TP53 (\texttt{ENSG00000141510}) and CENPA (\texttt{ENSG00000115163}) are always
retained.
After standardization, 40 principal components are computed from HVG expression.

\subsection{UMAP Visualization}

Two-dimensional UMAP coordinates are computed from the 40-component PCA embedding
using the \texttt{umap-learn} library \citep{mcinnes2018umap} with
\texttt{n\_neighbors}$=15$ and \texttt{min\_dist}$=0.1$.
The \texttt{n\_neighbors} parameter matches the cell k-nearest-neighbor graph
parameter used in heterogeneous graph construction, ensuring that the global
topology captured by UMAP reflects the same cell-similarity scale encoded in
the graph edges.

\subsection{Baselines and Evaluation Protocol}

Two unsupervised baselines are evaluated: k-means on HVG log-expression and
k-means on PCA coordinates.
Cluster assignments are mapped to condition labels using the Hungarian algorithm
on the training split \citep{luecken2022benchmarking}.
Supervised baselines use logistic regression, random forest (250 trees), and
multilayer perceptron (MLP; hidden layers 96--48) classifiers on 40-dimensional
PCA features.

To obtain variance estimates, all supervised baselines are evaluated with
five-fold stratified cross-validation.
Results are reported as mean $\pm$ standard deviation of macro-F1 across folds.
The full dataset (9,201 cells) is partitioned by \texttt{StratifiedKFold} with
\texttt{shuffle=True} and \texttt{random\_state=42}; a fresh \texttt{StandardScaler}
is fit on each training fold to prevent information leakage.
A fixed 75\%/25\% stratified split is retained as the primary held-out test set
for direct comparison with the HGT and homogeneous GCN, which cannot be embedded
into sklearn cross-validation without custom batching.

To test whether performance differences between models are statistically
significant, McNemar's test \citep{mcnemar1947} is applied pairwise on the
fixed held-out test split.
McNemar's test is appropriate for paired comparisons on the same test set because
it conditions on the number of discordant predictions rather than absolute counts,
providing higher power than chi-square for binary error comparison.
When the sum of discordant counts ($b+c$) is below 25, the exact binomial
variant is used; otherwise the chi-square approximation with continuity correction
is applied.

Embedding quality is quantified using the silhouette score
\citep{rousseeuw1987silhouette} and Calinski-Harabasz index
\citep{calinski1974dendrite} computed on the six-condition labels.
For silhouette score, a random subsample of 3,000 cells manages $O(n^2)$ memory
requirements; the Calinski-Harabasz index uses all cells.
These metrics are computed for PCA coordinates, UMAP coordinates, and the
64-dimensional HGT cell embeddings.

A label-efficiency experiment is conducted by training logistic regression on
stratified subsamples of 10\%--100\% of the fixed training split and evaluating
on the fixed test set.
This experiment characterizes how strongly graph-based neighborhood propagation
could benefit settings with limited label information.

An optional homogeneous GNN ablation uses only the cell--cell $k$-nearest-neighbor
graph, removing typed gene and module relations.
This ablation requires PyTorch Geometric and is run with
\texttt{src.train\_homogeneous\_gnn}.

\subsection{Heterogeneous Graph Construction}

The graph has three node types: cells, genes, and expression-derived gene modules.
Cell node features are the 40 PCA coordinates.
Gene node features summarize mean log-expression, variance, and detection frequency.
Module node features summarize the expression profiles of k-means gene clusters
(24 modules by default).

Four relation families define the graph: cell--gene expression, cell--cell
neighborhood, gene--gene coexpression, and gene--module membership.
Reciprocal edges are stored for directed HGT message passing so that all node
types can be updated.
Expression edges connect each cell to its top-40 expressed retained genes.
Cell-neighborhood edges are k-nearest-neighbor edges in PCA space ($k=15$).
Gene-coexpression edges connect each gene to its 8 nearest positively correlated
genes.
Gene-module edges connect each retained gene to its k-means-derived module.

\subsection{HGT Model}

The optional HGT model is implemented in PyTorch Geometric using \texttt{HGTConv}
\citep{pyg_hgtconv,pyg_heterodata}.
It projects each node type into a shared hidden dimension (64), applies two HGT
layers with 4 attention heads, and uses three classification heads on the final
cell embeddings: p53 status (2 classes), CENP-A status (3 classes), and the
six-condition joint label (6 classes).
The multitask objective is a weighted sum of cross-entropy losses
(weights: p53 $1.0$, CENP-A $1.0$, condition $1.25$).
The model is trained for 80 epochs with AdamW
($\mathrm{lr}=0.003$, $\mathrm{weight\_decay}=10^{-4}$).

\subsection{External Annotation}

External annotation is intentionally separated from the default pipeline.
The command \texttt{python3 -m src.external\_annotation} can query Ensembl REST
for gene symbols and g:Profiler for GO biological process and Reactome enrichment.
This step is not run automatically because it exports the candidate gene list to
third-party services.
